\summarychapter{2}{Wineskins: Old and New}{The Gospel to the Gentiles}

\begin{inbrief}
Within forty years a small Jewish sect in Jerusalem became a fellowship of many peoples spread across
the Roman Empire. The chapter traces that change from Stephen's death through Pentecost, the Greek-speaking
Jewish believers and the missions of Paul to the fall of Jerusalem in AD 70, when the ``old wineskin''
of Judaism finally tore.
\question{How did a Jewish Messiah, preaching a Jewish theme to a Jewish following, become the Savior of
people everywhere?}
\end{inbrief}

The Sanhedrin had failed to silence the ``Nazarenes,'' who even accused it of killing the Messiah. Then
Stephen went further and openly attacked the law of Moses and the temple. In his defense he retold
Israel's history to show that God may be worshiped apart from the temple, that Moses had foretold the
Messiah, and that his judges, like their fathers, resisted the Holy Spirit and had murdered the
Righteous One. The furious council dragged him outside the walls and stoned him.

\ssection{Christianity and Judaism}
For Shelley the death of Stephen, the first Christian martyr, holds the answer to the chapter's
question. The quarrel was about how to read the Old Testament. The scribes and Pharisees saw in it
God's law for his special people, covering every part of life. Meeting Jesus made the first believers
read it again, and they found a message for the whole world---God's promise to Abraham that all
peoples would be blessed in him. Stephen said plainly that law and temple were temporary: they pointed
beyond themselves to the Messiah who would fulfill righteousness for all people, and that Messiah was
Jesus. The proof lay in what followed the crucifixion. The apostles had gone into hiding with their
hopes destroyed; then women reported seeing Jesus alive, the tomb was found empty, and appearances of
the risen Jesus convinced even the skeptics. In Galilee he told them to wait in Jerusalem for the
baptism of the Holy Spirit.

\begin{sidebox}{Resurrection Faith}
Critics say that grieving, hopeful disciples produced visions of Jesus, even group hallucinations. The
first-century evidence points the other way. The disciples were defeated and took Jesus for one more
messiah executed by Rome; not even the empty tomb changed their minds. The encounters with the risen
Jesus created resurrection faith, not the reverse.
\end{sidebox}

\ssection{Pentecost}
Seven weeks after the crucifixion, at the feast of Pentecost, about 120 disciples were meeting in a
house when the Spirit came on them, with a sound like wind and what looked like tongues of fire. They
poured into the streets, and visitors to the city heard them in their own languages. At the temple
Peter explained the event as the outpouring of the Spirit promised by Joel for the ``last days'' and
traced it to Jesus: God had raised the crucified one and made him Lord and Messiah, as the Psalms had
said, ``and we are all witnesses.'' He called the crowd to repent and be baptized in the name of Jesus
for the forgiveness of sins and the gift of the Spirit. About three thousand accepted. That, says
Shelley, is how the church began, and Christians have held ever since that the cross, the resurrection
and the mission of the Spirit are the foundations of the faith; without the resurrection there is no
gospel, no salvation and no church. Stephen was the first to grasp that this could not be kept inside
the Pharisees' law. Jesus himself had hinted at the breach when he said that new wine is not poured
into old wineskins; the tearing of the old wineskins was the most important development of the first
century.

\begin{sidebox}{Tongue Speaking or Glossolalia}
Christians disagree about speaking in tongues. Some hold that the Spirit enables a person to give God's
message in a real language never learned. Others hold that it is ecstatic speech, not a human language;
in that case Acts describes a miracle of hearing, not of speaking: listeners understood in their own
language.
\end{sidebox}

\ssection{The First Community}
The first believers were Jews: Mary and other relatives of Jesus and the apostles, with Matthias chosen
to replace Judas Iscariot, who had killed himself. They kept the Jewish law, worshiped in synagogue and
temple, and looked like any other Jewish sect. They called their movement ``the Way''; soon they used
an Old Testament word for the assembly of Israel, in Greek \emph{ekklesia}, ``church,'' the gathering
of God's people. The temple authorities arrested the apostles and forbade them to preach the
resurrection, without effect; within two years the believers numbered several thousand. Two ceremonies
kept the death and resurrection of Jesus at the center. Baptism, unlike John's baptism for a kingdom
still to come, marked entry into a spiritual kingdom already announced though not yet fully revealed;
it expressed faith in Jesus as Lord. The Lord's Supper looked back to his betrayal and death as the
new covenant promised by Jeremiah and renewed the covenant of the believers with God and one another.

\ssection{The Hellenists}
Growth brought tension. More and more converts were Hellenists, Jews from all over the empire who had
settled in Jerusalem, spoke Greek, read the Septuagint (the Greek Old Testament) and were more open to
Gentiles and new ideas. When they complained that their widows were neglected in the daily
distribution, the apostles appointed seven Hellenists, among them Stephen and Philip, to oversee
it---perhaps the first \emph{deacons} (servants). Stephen's preaching in the Hellenist synagogues led
to the riot and his death, and vigilantes, among them the Pharisee Saul of Tarsus, began to seize
Nazarenes. This first bloodshed, about AD 36, widened the gap with Judaism and turned the faith into a
missionary movement. The Hellenist believers fled and founded churches in Samaria, Damascus, Antioch,
Tarsus, Cyprus and Egypt. Jerusalem sent Peter and John to Samaria and Barnabas of Cyprus to Antioch,
where unnamed men of Cyprus and Cyrene had taken the revolutionary step of evangelizing Gentiles.
Antioch, capital of Syria and third city of the empire with half a million people, was where the
disciples were first called ``Christians,'' at first a term of scorn. It replaced Jerusalem as the
center of mission, largely through Saul, who joined Barnabas there about AD 44.

\ssection{The Apostle Paul}
After Jesus, no one shaped Christianity more than Saul, or Paul. As a persecutor he could not believe in
a crucified, and therefore cursed, Messiah. Outside Damascus a blinding light and a voice---``Saul,
Saul, why are you persecuting me?''---overturned him, and Stephen's argument fell into place: the law
curses all who fail to keep it, but Christ redeemed us by becoming a curse for us. The law was given
for a time to show people that they cannot do God's will and to leave them only the gospel. The
persecutor became the persecuted. He was a man of three worlds: trained in the strictest Judaism under
Gamaliel, fluent in Greek language and thought, and a Roman citizen with freedom to travel.

His converts in Asia Minor and Greece were mostly pagans with sordid pasts, and the question of how to
teach them Christian morals divided the church. Palestinian Christians wanted Gentiles to keep the
Jewish law as well as believe. Paul answered from experience that righteousness by effort ends in
failure; people are accepted only by God's undeserved mercy, by grace, which comes from Christ's life,
death and resurrection. To the fear that justification by faith alone would breed loose living he
replied that whoever truly accepts Christ accepts his way and mind and will choose God's will. Shelley
notes that this dispute between the legally minded and the champions of grace has never left
Christianity.

Paul's first journey took him to Cyprus and Galatia; the second through Asia Minor to Troas and across
to Europe: Philippi, Thessalonica, Berea, Athens and Corinth, where he stayed a year and a half; the
third to Ephesus for more than two years. The pagan world was religious, and the mystery cults of dying
and rising gods (Isis, Mithras and others), with their ideas of immortality and of a struggle between
good and evil, gave him points of contact. Arrested in Jerusalem, he spent two years in custody at
Caesarea, appealed to the emperor and came to Rome, where he preached while awaiting trial. After
Nero's persecution in AD 64 nothing more is heard of him, though tradition speaks of a journey to
Spain. By then Gentile believers were uncircumcised, ignored the food laws and worshiped on the first
day of the week, the day of the resurrection.

\ssection{The Decline of Jerusalem}
The Jerusalem church stayed strictly Jewish and still suffered. About AD 41 Herod Agrippa I had James
the son of Zebedee killed; John may then have left the city; Peter was arrested, escaped, traveled
widely and at last died with Paul in Nero's persecution at Rome. James, ``the brother of the Lord,'' a
devout keeper of the law, led the church until the high priest had him murdered in AD 62, which left
it leaderless. Tension with Rome grew, sharpened by unemployment when the temple was finished in
AD 64, and in 66 the Jews revolted. After four years of war the forces of the emperor Vespasian, led
by Titus, took Jerusalem in AD 70, burned the temple and destroyed the city; the synagogues of
Palestine were burned in reprisal. The Jerusalem Christians, warned in a vision, had fled at the
outbreak. Pious Jews called this treason, and when Christian Jews were later barred from the
synagogues the break was complete. Christianity had become a Gentile movement; the old wineskin was
torn beyond repair. For practical purposes AD 70 ends the apostolic age. Most of the apostles were
dead, but their message would outlast persecution and become the dominant faith of the empire.

\begin{keyfacts}
\keypeople{Stephen; Peter; John; Matthias; Philip; Barnabas; Saul of Tarsus (Paul); Gamaliel; James the
son of Zebedee; James, ``the brother of the Lord''; Herod Agrippa I; Nero; Vespasian; Titus.}
\keydates{%
  \kd{c. AD 36}{Death of Stephen; the Hellenist believers scatter and begin missions}
  \kd{c. AD 41}{James the son of Zebedee killed by Herod Agrippa I}
  \kd{c. AD 44}{Saul joins Barnabas at Antioch}
  \kd{AD 62}{James, the brother of the Lord, murdered in Jerusalem}
  \kd{AD 64}{Nero's persecution of Christians in Rome}
  \kd{AD 66}{Jewish revolt against Rome begins}
  \kd{AD 70}{Titus destroys Jerusalem and the temple; end of the apostolic age}}
\keyterms{Pentecost; glossolalia; the Way; \emph{ekklesia}; baptism; Lord's Supper; Hellenists;
Septuagint; deacons; ``Christians''; grace; justification by faith; mystery cults.}
\end{keyfacts}
